\section{Introduction}
\label{sec:introduction}

Time-series event detection is the interface between long recordings and discrete decisions.
Clinical recordings, activity traces, industrial logs, audio streams, and scientific measurements can contain thousands or millions of samples, but the action taken downstream is often triggered by a much smaller set of times: an alarm, a changepoint, the beginning or end of an episode, or the next repeated action \citep{guralnik1999,aminikhanghahi2017,shoeb2010,goldberger2000,alexander2017,esper2023,azib2023universal}.
Those detected times decide which segment is reviewed, which alert is raised, which episode enters a database, or which candidate an expert must confirm, so a detector can be accurate over most samples and still fail by being late, early, duplicated, or poorly ranked at the few times that matter.
Modern event-detection benchmarks encode this decision interface directly.
They sort detections by confidence, match them to annotations inside task-specific tolerance windows, and summarize the resulting precision--recall curve.
The model is therefore judged not only by labels assigned to samples, but by how well it ranks the temporal decisions that the benchmark and downstream system will use.

A common reduction is to convert sparse event annotations into a samplewise segmentation problem.
Interval labels are expanded into per-sample states, a sequence model predicts those states, and event scores are recovered afterward from thresholds, transitions, or edge evidence.
This reduction is natural when deployment consumes calibrated state probabilities at every sample.
It becomes indirect when evaluation ignores most samples and rewards the timing of localized detections.
The issue is not the use of a time-discretized sequence output, which is often the practical interface for long recordings.
The issue is what that output is trained to estimate.
The same distinction appears whenever long context supports a short decision, repeated events can occur within one window, several event types are scored separately, or annotation uncertainty is mostly about time rather than class identity.
Sleep-event detection makes the contrast easy to see: a sleep interval can last for hours, while the benchmark gives credit for onset and wake-up predictions inside minute-scale windows.
The interval interior provides evidence for the events, but event AP matches the endpoints.

The broader detection literature already shows that supervision need not be uniform over an annotated region.
Pose, audio, and temporal-action systems often emphasize keypoints, sound events, starts, or ends rather than treating every background and interior sample symmetrically \citep{newell2016stacked,xiao2018simple,luo2021,yu2022,phan2015,gao2017turn,lin2018bsn,lin2019bmn}.
Recent time-series event work makes the same move away from hard per-sample classification by using overlap- or proximity-derived labels near annotated events \citep{azib2023universal}.
These precedents motivate localized supervision, but they do not by themselves define what that training signal should mean after smoothing, clipping, temporal downsampling, or nearby annotations.
When the final score is event AP, the invariant object is already present in the annotation: a time region contains zero, one, or several events.
A training signal for this setting should preserve those event occurrences while allowing the detector to choose a resolution, a timing-uncertainty model, and a decoder.

We introduce \emph{Boundary Density Likelihood} (BDL), a likelihood objective for sequence event detectors.
BDL treats each annotation as one event occurrence of its type: smoothing can redistribute that occurrence to represent timing uncertainty, temporal binning sums it into model outputs, and a Poisson negative log-likelihood fits nonnegative predictions with units of expected events per bin.
The units are therefore fixed by the annotation, unlike normalized position distributions, unscaled proximity scores, or samplewise state labels.
This construction applies to empty windows, repeated events, point events, onset-only labels, interval endpoints, and multi-type streams without requiring paired endpoints.

This paper makes three contributions.
First, it gives sequence detectors a training signal aligned with the objects matched by event-level average precision: how many scored events fall in each output bin, rather than only which state label holds at each sample.
Second, it fixes the units of localized event supervision: smoothing, clipping, and temporal binning can change where an annotation is placed on the output timeline, but not how many annotated events it represents.
Third, it evaluates the resulting likelihood objective in controlled detector comparisons rather than as a standalone architecture change.
On sleep-event detection, changing the supervision improves strict event AP across recurrent and convolutional backbones, while class weighting and focal segmentation do not recover the same localization performance.
CHB-MIT, point-event, kernel, architecture, streaming-context, and state-reconstruction ablations then separate the effect of likelihood supervision from representation, output stride, calibration, event structure, and decoding choices.
